\documentclass[11pt]{amsart}

\usepackage[margin=1in]{geometry}
\usepackage{amsmath,amssymb,amsthm}
\usepackage{graphicx}
\usepackage[hidelinks]{hyperref}

\newtheorem{theorem}{Theorem}
\newtheorem{corollary}[theorem]{Corollary}
\newtheorem{lemma}[theorem]{Lemma}

\newcommand{\N}{\mathbb N}
\newcommand{\I}{\mathcal I}
\newcommand{\B}{\mathcal B}
\newcommand{\Sum}{\operatorname{Sum}}
\newcommand{\ellq}{\ell_q}

\title{A positive answer to the cotype summability representation question}
\author{Agent lane 19}
\date{2026-07-03}

\begin{document}
\maketitle

\section*{Source problem}

The source paper is Lopez-Abad, Olmos-Prieto, and Uzcategui-Aylwin,
\emph{\texorpdfstring{$F_\sigma$}{F-sigma}-ideals, colorings, and representation
in Banach spaces}, arXiv:2501.15643.  Question 3.9 asks whether an ideal that is
\(\B\)-representable on a Banach space of non-trivial cotype \(q\) can be
\(\B\)-represented, still in cotype \(q\), by a sequence \((x_n)\) for which
\(\sum_n \|x_n\|^q=\infty\).

\begin{center}
\includegraphics[width=\textwidth]{figures/open_problem_crop.png}
\end{center}

Here ideals are understood in the convention used by the source and by
Filipow--Tryba: they are proper ideals on \(\N\) containing all finite subsets.
The improper family \(\mathcal P(\N)\) is therefore outside the convention.

\section*{Definitions}

For a sequence \(\mathbf x=(x_n)\) in a Banach space \(X\), write
\[
  \B_X(\mathbf x)
  = \left\{A\subseteq\N:
      \sup_{F\subseteq A,\ F\text{ finite}}
      \left\|\sum_{n\in F} x_n\right\|_X <\infty
    \right\}.
\]
For weights \(w=(w_n)\) with \(0\leq w_n<\infty\), write
\[
  \Sum(w)=\left\{A\subseteq\N:\sum_{n\in A}w_n<\infty\right\}.
\]
This summable ideal is non-trivial precisely when \(\sum_n w_n=\infty\).

\section*{Proof Intuition}

The source paper proves the one-sided cotype obstruction:
if \(X\) has cotype \(q\), then every \(\B_X(\mathbf x)\)-set has finite
\(\sum \|x_n\|^q\).  Thus a representation with divergent total
\(\sum_n\|x_n\|^q\) can exist only when the ideal sits inside some non-trivial
summable ideal.

The converse is the useful trick.  If
\(\I=\B_X(\mathbf x)\subseteq \Sum(w)\) and \(\sum_n w_n=\infty\), add one
orthogonal \(\ell_q\)-coordinate with length \(w_n^{1/q}\) to each vector:
\[
  y_n=(x_n,w_n^{1/q}e_n)\in X\oplus_q\ell_q.
\]
Finite sums then split exactly into the original \(X\)-partial sum and the
weight sum.  Hence the new bounded-partial-sum ideal is
\(\B_X(\mathbf x)\cap\Sum(w)=\I\), while the total \(q\)-mass diverges because of
the added \(\ell_q\)-coordinate.

\section*{Main theorem}

\begin{theorem}
\label{thm:equivalence}
Let \(2\le q<\infty\), and let
\(\I=\B_X((x_n))\) for a sequence in a Banach space \(X\) of cotype \(q\).
Then the following are equivalent.
\begin{enumerate}
\item There are a Banach space \(Y\) of cotype \(q\) and a sequence
  \((y_n)\subset Y\) such that
  \[
    \I=\B_Y((y_n))
    \quad\text{and}\quad
    \sum_{n=1}^\infty \|y_n\|_Y^q=\infty .
  \]
\item \(\I\) is contained in a non-trivial summable ideal.
\end{enumerate}
\end{theorem}

\begin{proof}
Assume first that (1) holds.  Proposition 3.7 of the source paper says that if
\(Y\) has cotype \(q\), then
\[
  \B_Y((y_n))\subseteq \Sum((\|y_n\|_Y^q)_n).
\]
Since \(\sum_n\|y_n\|_Y^q=\infty\), the ideal
\(\Sum((\|y_n\|_Y^q)_n)\) is non-trivial.  Thus \(\I\) is contained in a
non-trivial summable ideal.

Conversely assume that \(\I\subseteq\Sum(w)\), where
\(w_n\ge 0\) and \(\sum_n w_n=\infty\).  Let
\[
  Y=X\oplus_q\ell_q,
  \qquad
  \|(u,a)\|_Y=(\|u\|_X^q+\|a\|_{\ell_q}^q)^{1/q},
\]
and define \(y_n=(x_n,w_n^{1/q}e_n)\).

The space \(Y\) has cotype \(q\).  Indeed, if \(X\) has cotype \(q\) with
constant \(C_X\), and \(\ell_q\) has cotype \(q\) with constant \(C_q\), then
for every finite family \(z_j=(u_j,v_j)\in Y\),
\[
\begin{aligned}
\sum_j \|z_j\|_Y^q
 &= \sum_j\|u_j\|_X^q+\sum_j\|v_j\|_{\ell_q}^q  \\
 &\le C_X\,\mathbb E_\varepsilon
       \left\|\sum_j\varepsilon_j u_j\right\|_X^q
     + C_q\,\mathbb E_\varepsilon
       \left\|\sum_j\varepsilon_j v_j\right\|_{\ell_q}^q  \\
 &\le \max(C_X,C_q)\,\mathbb E_\varepsilon
       \left\|\sum_j\varepsilon_j z_j\right\|_Y^q .
\end{aligned}
\]

For every finite \(F\subseteq\N\), the definition of the \(\ell_q\)-sum gives
\[
  \left\|\sum_{n\in F}y_n\right\|_Y^q
  =
  \left\|\sum_{n\in F}x_n\right\|_X^q
  +\sum_{n\in F}w_n .
\]
Therefore, for every \(A\subseteq\N\),
\[
  A\in \B_Y((y_n))
  \Longleftrightarrow
  A\in\B_X((x_n))\text{ and }A\in\Sum(w).
\]
Since \(\I=\B_X((x_n))\subseteq\Sum(w)\), this yields
\(\B_Y((y_n))=\I\).  Finally,
\[
  \sum_n\|y_n\|_Y^q
  =
  \sum_n(\|x_n\|_X^q+w_n)
  \ge \sum_n w_n
  =\infty .
\]
This proves (1).
\end{proof}

\begin{corollary}
\label{cor:question}
Question 3.9 of arXiv:2501.15643 has an affirmative answer for every proper
\(\B\)-ideal that is representable on a space of cotype \(q\).
\end{corollary}

\begin{proof}
The source paper recalls that \(\B\)-ideals are exactly the nonpathological
\(F_\sigma\) ideals.  Filipow--Tryba, \emph{Path of pathology},
Theorem 11.2, proves that every nonpathological \(F_\sigma\) ideal is an
intersection of summable ideals; in particular it is contained in a non-trivial
summable ideal.  Apply Theorem~\ref{thm:equivalence}.
\end{proof}

\section*{Verification notes}

The construction preserves the original cotype exponent \(q\), not merely some
larger finite cotype, because the added coordinate is \(\ell_q\) and the direct
sum is the \(\ell_q\)-sum.  The equality
\(\B_Y((y_n))=\B_X((x_n))\cap\Sum(w)\) is exact and uses only the finite-partial
sum characterization of weak unconditional convergence.

The only external nontrivial input beyond the source paper is Filipow--Tryba's
Theorem 11.2.  It supplies the non-trivial summable ideal needed in
Theorem~\ref{thm:equivalence}.  The theorem is stated for proper ideals, matching
the convention in the source paper.

\section*{Novelty and limitations}

I searched the run indexes for arXiv:2501.15643, the source title, cotype,
\(\B\)-representability, and the summable-representation question.  I also ran a
web search for exact phrases from Question 3.9 and for the source/supporting
paper titles.  The search found the source paper and Filipow--Tryba's supporting
paper, but no separate statement of the direct-sum cotype representation theorem
above.

The answer relies on the standard characterization of \(\B\)-ideals as
nonpathological \(F_\sigma\) ideals.  If a reader uses a nonstandard convention
allowing \(\mathcal P(\N)\) as an ideal, then the conclusion must exclude that
improper case: \(\mathcal P(\N)\) can be represented by the zero sequence, but it
cannot be represented in cotype \(q\) with divergent \(\sum\|x_n\|^q\), because
Proposition 3.7 would force \(\N\in \Sum((\|x_n\|^q))\).

\section*{Human review recommendation}

This is a short candidate full solution.  The main review points are:
verify the convention that all ideals in the source question are proper; verify
the cited equivalence between \(\B\)-ideals and nonpathological \(F_\sigma\)
ideals; and check that the \(\ell_q\)-sum cotype estimate is acceptable with the
source paper's normalization of cotype constants.

\begin{thebibliography}{9}

\bibitem{LOU}
J. Lopez-Abad, V. Olmos-Prieto, and C. Uzcategui-Aylwin,
\emph{\(F_\sigma\)-ideals, colorings, and representation in Banach spaces},
arXiv:2501.15643.

\bibitem{FT}
R. Filipow and J. Tryba,
\emph{Path of pathology},
arXiv:2501.00503.

\end{thebibliography}

\end{document}
